% !TEX root = ../main.tex
%% Chapitre 9 — Calcul matriciel et systèmes linéaires
%% Programme officiel MPSI/MP2I (BO spécial n°1 du 11/02/2021), premier semestre,
%% section « Calcul matriciel et systèmes linéaires » : a) opérations sur les
%% matrices ; b) opérations élémentaires ; c) systèmes linéaires ; anneau des
%% matrices carrées ; matrices inversibles, GL_n(K), calcul d'inverse,
%% matrices triangulaires.
%% Résultats vérifiés par verifications/chapitre_09.py.

\begin{ouverture}[matrices]{9}{Calcul matriciel et systèmes linéaires}{Partie 1 · Premier semestre}
  {Produit matriciel\enspace·\enspace Transposée\enspace·\enspace Anneau $\mathcal{M}_n(\K)$\enspace·\enspace
   Inversibilité\enspace·\enspace Systèmes $AX=B$}
  \objectif{Calculer avec les matrices : combinaisons linéaires, produit, transposée.}
  \objectif{Manipuler les matrices élémentaires et le symbole de Kronecker.}
  \objectif{Calculer dans l'anneau $\mathcal{M}_n(\K)$ : puissances, binôme, nilpotents.}
  \objectif{Reconnaître une matrice inversible et calculer son inverse.}
  \objectif{Exploiter les matrices triangulaires et diagonales.}
  \objectif{Traduire un système linéaire sous la forme $AX=B$ et structurer ses solutions.}
  \prerequis{Systèmes linéaires et algorithme du pivot (chapitre 2).}
  \prerequis{Anneaux, groupes, éléments inversibles et nilpotents (chapitre 8).}
  \prerequis{Formule du binôme (chapitre 2) et trigonométrie (chapitre 2).}
\end{ouverture}

%% =====================================================================
%%  L'ESSENTIEL
%% =====================================================================
\essentiel

\essSection{1}{L'espace des matrices}

\begin{definition}[Matrices]
Pour $n,p\geq1$ et $\K=\R$ ou $\Cx$, on note $\mathcal{M}_{n,p}(\K)$
l'ensemble des \textbf{matrices} à $n$ lignes et $p$ colonnes à
coefficients dans $\K$ ; on écrit $A=(a_{i,j})$ où $a_{i,j}$ est le
coefficient en ligne $i$, colonne $j$. L'addition et la multiplication par
un scalaire se font coefficient par coefficient. La \textbf{matrice
élémentaire} $E_{i,j}$ a un unique coefficient non nul, un $1$ en position
$(i,j)$ : toute matrice est combinaison linéaire de matrices élémentaires,
$A=\sum_{i,j}a_{i,j}E_{i,j}$. On utilise le \textbf{symbole de Kronecker} :
$\delta_{i,j}=1$ si $i=j$, $0$ sinon.
\end{definition}

\essSection{2}{Le produit matriciel}

\noindent
\begin{minipage}[c]{0.56\linewidth}
\begin{definition}
Pour $A\in\mathcal{M}_{n,p}(\K)$ et $B\in\mathcal{M}_{p,q}(\K)$, le
\textbf{produit} $AB\in\mathcal{M}_{n,q}(\K)$ a pour coefficients
\[
  (AB)_{i,j}=\sum_{k=1}^{p}a_{i,k}\,b_{k,j} :
\]
la \emph{ligne $i$ de $A$} rencontre la \emph{colonne $j$ de $B$}. Le
produit est \emph{bilinéaire} et \emph{associatif}. Si $X$ est une matrice
colonne, $AX$ est la combinaison linéaire des colonnes de $A$ dont les
coefficients sont ceux de $X$.
\end{definition}
\end{minipage}\hfill
\begin{minipage}[c]{0.4\linewidth}
\centering
\begin{tikzpicture}[scale=0.44,line cap=round]
  % B (en haut à droite)
  \draw[otMuted,line width=0.5pt,rounded corners=0.8mm] (4.6,4.6) rectangle (9.4,9.4);
  \fill[otGreenTint] (6.8,4.75) rectangle (7.9,9.25);
  \foreach \i in {0,...,3}{\foreach \j in {0,...,3}{
    \fill[otMuted!55] (5.3+\j*1.1,5.3+\i*1.1) circle (0.16);}}
  \foreach \i in {0,...,3}{\fill[otGreenDark] (7.5,5.3+\i*1.1) circle (0.19);}
  \node[font=\scriptsize,text=otMuted] at (9.9,9.0) {$B$};
  % A (en bas à gauche)
  \draw[otMuted,line width=0.5pt,rounded corners=0.8mm] (-0.4,-0.4) rectangle (4.4,4.4);
  \fill[otGreenTint] (-0.25,1.65) rectangle (4.25,2.75);
  \foreach \i in {0,...,3}{\foreach \j in {0,...,3}{
    \fill[otMuted!55] (0.3+\j*1.1,0.3+\i*1.1) circle (0.16);}}
  \foreach \j in {0,...,3}{\fill[otNavy] (0.3+\j*1.1,2.5) circle (0.19);}
  \node[font=\scriptsize,text=otMuted] at (-1.0,0.2) {$A$};
  % AB (en bas à droite)
  \draw[otMuted,line width=0.5pt,rounded corners=0.8mm] (4.6,-0.4) rectangle (9.4,4.4);
  \fill[otGreen] (7.5,2.5) circle (0.3);
  \node[font=\scriptsize,text=otMuted] at (9.9,0.2) {$AB$};
  \draw[otRule!90!otMuted,densely dotted,line width=0.5pt] (4.4,2.5) -- (7.1,2.5);
  \draw[otRule!90!otMuted,densely dotted,line width=0.5pt] (7.5,4.6) -- (7.5,2.9);
\end{tikzpicture}\\[0.5mm]
\legende{Le coefficient $(i,j)$ de $AB$ : la ligne $i$ de $A$ contre la colonne $j$ de $B$.}
\end{minipage}

\begin{propriete}[Produit des matrices élémentaires]
Dans $\mathcal{M}_n(\K)$ :
$E_{i,j}\,E_{k,l}=\delta_{j,k}\,E_{i,l}$ — le produit est nul sauf si les
indices « se raccordent ». Première conséquence : le produit matriciel
n'est \emph{pas} commutatif, et un produit peut être nul sans qu'aucun
facteur ne le soit ($E_{1,2}E_{1,2}=0$).
\end{propriete}

\essSection{3}{Transposée, matrices symétriques}

\begin{definition}
La \textbf{transposée} $A^{\top}\in\mathcal{M}_{p,n}(\K)$ de
$A\in\mathcal{M}_{n,p}(\K)$ échange lignes et colonnes :
$(A^{\top})_{i,j}=a_{j,i}$. La transposition est linéaire, involutive, et
\emph{renverse les produits} :
\[
  (AB)^{\top}=B^{\top}A^{\top}.
\]
Une matrice carrée est \textbf{symétrique} si $A^{\top}=A$ (ensemble
$\mathcal{S}_n(\K)$), \textbf{antisymétrique} si $A^{\top}=-A$ (ensemble
$\mathcal{A}_n(\K)$). Toute matrice carrée se décompose de manière unique
en somme d'une symétrique et d'une antisymétrique :
$A=\frac{A+A^{\top}}{2}+\frac{A-A^{\top}}{2}$.
\end{definition}

\essSection{4}{L'anneau $\mathcal{M}_n(\K)$}

\begin{theoreme}
Muni de l'addition et du produit, $\mathcal{M}_n(\K)$ est un \emph{anneau},
d'unité la matrice identité $I_n$ ; les matrices \emph{scalaires}
$\lambda I_n$ commutent avec tout le monde. Pour $n\geq2$, cet anneau est
\textbf{non commutatif} et \textbf{non intègre} : il possède des diviseurs
de zéro et des éléments \emph{nilpotents}, comme
$N=E_{1,2}$ ($N^{2}=0$).
\end{theoreme}

\begin{savoirfaire}[Puissances d'une matrice]
\begin{itemize}[label={\color{otGreenDark}\textbullet},leftmargin=1.4em,itemsep=0.6mm,topsep=0.5mm]
  \item \etape{Diagonale} $\operatorname{diag}(d_1,\dots,d_n)^{k}
    =\operatorname{diag}(d_1^{k},\dots,d_n^{k})$ ; plus généralement, le
    produit de matrices diagonales (resp. triangulaires supérieures) est
    diagonal (resp. triangulaire supérieur).
  \item \etape{Binôme} si $AB=BA$, alors
    $(A+B)^{k}=\sum\binom{k}{j}A^{j}B^{k-j}$. Cas d'école : $A=\lambda I+N$
    avec $N$ nilpotente — la somme est alors \emph{finie}.
  \item \etape{Récurrence} conjecturer la forme de $A^{k}$ sur les
    premières puissances, puis démontrer par récurrence.
\end{itemize}
\end{savoirfaire}

\essSection{5}{Matrices inversibles}

\begin{definition}[Groupe linéaire]
$A\in\mathcal{M}_n(\K)$ est \textbf{inversible} s'il existe $B$ telle que
$AB=BA=I_n$ ; l'inverse $A^{-1}$ est alors unique. Les matrices
inversibles forment un groupe, le \textbf{groupe linéaire}
$\mathrm{GL}_n(\K)$, avec
\[
  (AB)^{-1}=B^{-1}A^{-1}
  \qquad\text{et}\qquad
  (A^{\top})^{-1}=(A^{-1})^{\top}.
\]
\end{definition}

\begin{exemple}[Le cas $2\times2$]
$A=\begin{pmatrix}a&b\\c&d\end{pmatrix}$ est inversible si et seulement si
$ad-bc\neq0$, et alors
\[
  A^{-1}=\frac{1}{ad-bc}\begin{pmatrix}d&-b\\-c&a\end{pmatrix}.
\]
\end{exemple}

\begin{savoirfaire}[Calculer un inverse]
\begin{itemize}[label={\color{otGreenDark}\textbullet},leftmargin=1.4em,itemsep=0.6mm,topsep=0.5mm]
  \item \etape{Par le système} résoudre $AX=Y$ d'inconnue $X$ : si la
    solution s'écrit $X=BY$ pour toute colonne $Y$, alors $A$ est
    inversible et $A^{-1}=B$.
  \item \etape{Par les opérations élémentaires} chaque opération
    élémentaire sur les lignes est la multiplication à gauche par une
    matrice inversible ; en transformant $A$ en $I_n$, les mêmes
    opérations transforment $I_n$ en $A^{-1}$ (disposition
    $(A\mid I_n)\to(I_n\mid A^{-1})$).
\end{itemize}
Les opérations élémentaires \emph{préservent l'inversibilité}.
\end{savoirfaire}

\begin{theoreme}[Matrices triangulaires]
Une matrice triangulaire est inversible si et seulement si ses coefficients
diagonaux sont tous non nuls ; son inverse est alors triangulaire de même
type. En particulier,
$\operatorname{diag}(d_1,\dots,d_n)^{-1}
=\operatorname{diag}(d_1^{-1},\dots,d_n^{-1})$ lorsque tous les $d_i$ sont
non nuls.
\end{theoreme}

\essSection{6}{Systèmes linéaires : la vision matricielle}

\begin{definition}
Tout système linéaire s'écrit $AX=B$, avec $A\in\mathcal{M}_{n,p}(\K)$ la
matrice des coefficients, $X$ la colonne des inconnues et $B$ le second
membre ; le système \textbf{homogène} associé est $AX=0$. Le système est
\textbf{compatible} si et seulement si $B$ est combinaison linéaire des
colonnes de $A$.
\end{definition}

\begin{theoreme}[Structure des solutions]
Si $AX=B$ est compatible, de solution particulière $X_0$, alors ses
solutions sont exactement les
\[
  X=X_0+Y,\qquad\text{où $Y$ parcourt les solutions de }AX=0 :
\]
« solution particulière $+$ solutions de l'homogène » — la même structure
que pour les équations différentielles linéaires (chapitre 4). L'algorithme
du pivot (chapitre 2) se relit comme une suite d'opérations élémentaires
sur les lignes de $(A\mid B)$. Enfin, si $A$ est \emph{carrée inversible},
le système possède l'unique solution $X=A^{-1}B$.
\end{theoreme}

\begin{attention}[Pièges fréquents]
\begin{itemize}[label={\color{otOrange}\textbullet},leftmargin=1.4em,itemsep=0.6mm,topsep=0.5mm]
  \item $AB\neq BA$ en général : développer $(A+B)^{2}$ donne
    $A^{2}+AB+BA+B^{2}$, et le binôme exige $AB=BA$.
  \item $AB=0$ n'implique pas $A=0$ ou $B=0$ : on ne simplifie pas par une
    matrice non inversible.
  \item Les tailles doivent s'enchaîner : $AB$ exige
    (colonnes de $A$) $=$ (lignes de $B$).
  \item $(AB)^{\top}=B^{\top}A^{\top}$ et $(AB)^{-1}=B^{-1}A^{-1}$ :
    l'ordre se renverse dans les deux cas.
  \item Une matrice à coefficients tous non nuls peut être non inversible,
    et une triangulaire à diagonale sans zéro l'est toujours : c'est la
    \emph{diagonale} qui décide pour les triangulaires.
  \item Un système compatible à matrice non inversible a une infinité de
    solutions : toujours donner la \emph{structure} ($X_0$ + homogène).
\end{itemize}
\end{attention}

%% =====================================================================
%%  EXERCICES
%% =====================================================================
\exercices
\rubrique{Applications}

\begin{exercice}[1]{Combinaisons linéaires}
On pose $A=\begin{pmatrix}1&2\\3&-1\end{pmatrix}$ et
$B=\begin{pmatrix}0&1\\2&5\end{pmatrix}$.
\begin{questions}
  \item Calculer $A+B$, $2A-3B$ et $A-B$.
  \item Écrire $A$ comme combinaison linéaire des matrices élémentaires
    $E_{1,1}$, $E_{1,2}$, $E_{2,1}$, $E_{2,2}$.
\end{questions}
\end{exercice}

\begin{exercice}[1]{Produits matriciels}
Avec les matrices $A$ et $B$ de l'exercice 1, et
$C=\begin{pmatrix}1&0&2\\0&1&-1\end{pmatrix}$,
$D=\begin{pmatrix}1&1\\2&0\\0&3\end{pmatrix}$ :
\begin{questions}
  \item Calculer $AB$ et $BA$. Que constate-t-on ?
  \item Calculer $CD$. Le produit $DC$ a-t-il un sens ? De quelle taille ?
  \item Pour $X=\begin{pmatrix}2\\-1\end{pmatrix}$, calculer $AX$ et
    vérifier que c'est la combinaison correspondante des colonnes de $A$.
\end{questions}
\end{exercice}

\begin{exercice}[1]{Matrices élémentaires et symbole de Kronecker}
Dans $\mathcal{M}_3(\K)$ :
\begin{questions}
  \item Calculer $E_{1,2}E_{2,3}$, puis $E_{1,2}E_{1,2}$.
  \item Démontrer la formule générale
    $E_{i,j}E_{k,l}=\delta_{j,k}E_{i,l}$.
  \item En déduire deux matrices non nulles dont le produit est nul.
\end{questions}
\end{exercice}

\begin{exercice}[1]{Transposées, symétriques, antisymétriques}
\begin{questions}
  \item Vérifier sur les matrices $A$ et $B$ de l'exercice 1 que
    $(AB)^{\top}=B^{\top}A^{\top}$.
  \item Décomposer $A$ en somme d'une matrice symétrique et d'une matrice
    antisymétrique.
  \item Montrer que cette décomposition est unique pour toute matrice
    carrée.
\end{questions}
\end{exercice}

\begin{exercice}[1]{Premières puissances}
\begin{questions}
  \item Pour $D=\operatorname{diag}(2,-3)$, donner $D^{n}$.
  \item Pour $T=\begin{pmatrix}1&1\\0&1\end{pmatrix}$, calculer $T^{2}$,
    $T^{3}$, conjecturer $T^{n}$ et le démontrer par récurrence.
  \item Pour $J=\begin{pmatrix}1&1\\1&1\end{pmatrix}$, vérifier que
    $J^{2}=2J$ et en déduire $J^{n}$ pour $n\geq1$.
\end{questions}
\end{exercice}

\begin{exercice}[1]{Inversibles $2\times2$}
\begin{questions}
  \item Démontrer le critère du cours : $A=\begin{pmatrix}a&b\\c&d\end{pmatrix}$
    est inversible si et seulement si $ad-bc\neq0$, avec la formule de
    l'inverse. \emph{On pourra calculer
    $A\begin{pmatrix}d&-b\\-c&a\end{pmatrix}$.}
  \item Les matrices $\begin{pmatrix}2&5\\1&3\end{pmatrix}$ et
    $\begin{pmatrix}1&2\\2&4\end{pmatrix}$ sont-elles inversibles ? Donner
    l'inverse le cas échéant.
\end{questions}
\end{exercice}

\begin{exercice}[1]{Un système, vu matriciellement}
On considère le système
$\begin{cases}x+y+z=6\\x+2y+3z=14\\2x+y-z=1\end{cases}$ (chapitre 2).
\begin{questions}
  \item L'écrire sous la forme $AX=B$.
  \item Le résoudre, et en déduire que $B$ est combinaison linéaire des
    colonnes de $A$, avec les coefficients que l'on précisera.
\end{questions}
\end{exercice}

\begin{exercice}[1]{Un inverse par la résolution de $AX=Y$}
Soit $A=\begin{pmatrix}1&2\\1&3\end{pmatrix}$.
\begin{questions}
  \item Résoudre le système $AX=Y$, d'inconnue
    $X=\begin{pmatrix}x\\y\end{pmatrix}$, pour un second membre
    $Y=\begin{pmatrix}u\\v\end{pmatrix}$ quelconque.
  \item En déduire que $A$ est inversible et donner $A^{-1}$ ; vérifier
    par le produit.
\end{questions}
\end{exercice}

\rubrique{Entraînement}

\begin{exercice}[2]{Diviseurs de zéro et nilpotents}
\begin{questions}
  \item Trouver deux matrices $A,B\in\mathcal{M}_2(\K)$ non nulles telles
    que $AB=0$ mais $BA\neq0$.
  \item Montrer que $N=\begin{pmatrix}0&1\\0&0\end{pmatrix}$ est
    nilpotente, et calculer $(I_2+N)^{-1}$ à l'aide du résultat du
    chapitre 8 sur $1-x$ nilpotent.
  \item Montrer qu'une matrice nilpotente n'est jamais inversible.
\end{questions}
\end{exercice}

\begin{exercice}[2]{Le binôme matriciel}
Soit $N=\begin{pmatrix}0&1&0\\0&0&1\\0&0&0\end{pmatrix}$ et $A=I_3+N$.
\begin{questions}
  \item Calculer $N^{2}$ et $N^{3}$.
  \item Justifier que la formule du binôme s'applique à $A=I_3+N$, et en
    déduire $A^{n}$ pour tout $n\in\N$.
  \item Vérifier la formule obtenue pour $n=2$ et $n=3$.
\end{questions}
\end{exercice}

\begin{exercice}[2]{Puissances : deux méthodes}
Soit $A=\begin{pmatrix}2&1\\0&2\end{pmatrix}$.
\begin{questions}
  \item Écrire $A=2I_2+N$ avec $N$ nilpotente, et calculer $A^{n}$ par la
    formule du binôme.
  \item Retrouver le résultat par récurrence.
\end{questions}
\end{exercice}

\begin{exercice}[2]{Suites couplées}
On définit deux suites par $u_0=v_0=1$ et, pour tout $n$ :
$u_{n+1}=2u_n+v_n$, $v_{n+1}=2v_n$.
\begin{questions}
  \item Écrire la récurrence sous la forme $X_{n+1}=AX_n$, avec
    $X_n=\begin{pmatrix}u_n\\v_n\end{pmatrix}$, et justifier que
    $X_n=A^{n}X_0$.
  \item À l'aide de l'exercice précédent, exprimer $u_n$ et $v_n$ en
    fonction de $n$.
\end{questions}
\end{exercice}

\begin{exercice}[2]{Symétriques : produits et pièges}
\begin{questions}
  \item Montrer que, pour toute matrice $A\in\mathcal{M}_{n,p}(\K)$, la
    matrice $AA^{\top}$ est symétrique.
  \item Montrer qu'un produit de deux matrices symétriques n'est pas
    toujours symétrique, à l'aide d'un contre-exemple $2\times2$.
  \item Montrer que, pour $A$ et $B$ symétriques : $AB$ est symétrique si
    et seulement si $AB=BA$.
\end{questions}
\end{exercice}

\begin{exercice}[2]{Matrices qui commutent}
\begin{questions}
  \item Déterminer toutes les matrices de $\mathcal{M}_2(\K)$ qui
    commutent avec $D=\operatorname{diag}(1,2)$.
  \item Déterminer toutes les matrices qui commutent avec
    $N=\begin{pmatrix}0&1\\0&0\end{pmatrix}$.
  \item Montrer que l'ensemble des matrices commutant avec une matrice
    $M$ fixée est stable par combinaison linéaire et par produit.
\end{questions}
\end{exercice}

\begin{exercice}[2]{Gauss-Jordan : un inverse $3\times3$}
Soit $A=\begin{pmatrix}1&1&1\\0&1&2\\1&2&2\end{pmatrix}$.
\begin{questions}
  \item En appliquant des opérations élémentaires au tableau
    $(A\mid I_3)$, calculer $A^{-1}$.
  \item Vérifier le résultat en calculant $AA^{-1}$.
  \item En déduire la solution du système $AX=\begin{pmatrix}2\\1\\3\end{pmatrix}$.
\end{questions}
\end{exercice}

\begin{exercice}[2]{Matrices triangulaires}
\begin{questions}
  \item Montrer que le produit de deux matrices triangulaires supérieures
    de $\mathcal{M}_3(\K)$ est triangulaire supérieur, et que les
    coefficients diagonaux du produit sont les produits des coefficients
    diagonaux.
  \item Justifier qu'une matrice triangulaire est inversible si et
    seulement si ses coefficients diagonaux sont non nuls.
    \emph{Raisonner sur le système triangulaire $TX=Y$.}
  \item Calculer l'inverse de
    $T=\begin{pmatrix}1&2&3\\0&4&5\\0&0&6\end{pmatrix}$, et vérifier
    qu'il est triangulaire supérieur.
\end{questions}
\end{exercice}

\begin{exercice}[2]{Un paramètre dans le système}
Pour $m\in\R$, on pose $A_m=\begin{pmatrix}1&1\\1&m\end{pmatrix}$ et
$B=\begin{pmatrix}2\\3\end{pmatrix}$.
\begin{questions}
  \item Pour quelles valeurs de $m$ la matrice $A_m$ est-elle inversible ?
    Résoudre alors $A_mX=B$.
  \item Pour $m=1$ : le système $A_1X=B$ est-il compatible ? Et le système
    $A_1X=\begin{pmatrix}2\\2\end{pmatrix}$ ? Décrire dans ce dernier cas
    l'ensemble des solutions sous la forme « particulière $+$ homogène ».
\end{questions}
\end{exercice}

\begin{exercice}[2]{Matrices de rotation}
Pour $\theta\in\R$, on pose
$R(\theta)=\begin{pmatrix}\cos\theta&-\sin\theta\\\sin\theta&\cos\theta\end{pmatrix}$.
\begin{questions}
  \item Montrer que $R(a)R(b)=R(a+b)$ pour tous réels $a,b$.
  \item En déduire que $R(\theta)$ est inversible, avec
    $R(\theta)^{-1}=R(-\theta)=R(\theta)^{\top}$, et que
    $R(\theta)^{n}=R(n\theta)$.
  \item Montrer que $\theta\mapsto R(\theta)$ est un morphisme de groupes
    de $(\R,+)$ dans $\mathrm{GL}_2(\R)$, et déterminer son noyau.
\end{questions}
\end{exercice}

\rubrique{Approfondissement}

\begin{exercice}[3]{Le centre de $\mathcal{M}_n(\K)$}
On cherche les matrices $Z\in\mathcal{M}_n(\K)$ qui commutent avec
\emph{toutes} les matrices.
\begin{questions}
  \item Justifier qu'il suffit que $Z$ commute avec toutes les matrices
    élémentaires $E_{i,j}$.
  \item En examinant les coefficients de $ZE_{i,j}$ et de $E_{i,j}Z$,
    montrer que $Z$ est nécessairement scalaire : $Z=\lambda I_n$.
  \item Conclure : le centre de l'anneau $\mathcal{M}_n(\K)$ est
    l'ensemble des matrices scalaires.
\end{questions}
\end{exercice}

\begin{exercice}[3]{Une équation matricielle}
On cherche toutes les matrices $X\in\mathcal{M}_2(\R)$ telles que
$X^{2}=T$, où $T=\begin{pmatrix}1&1\\0&1\end{pmatrix}$.
\begin{questions}
  \item Montrer que toute solution $X$ commute avec $T$.
  \item En déduire, à l'aide de l'exercice 14, que $X$ est de la forme
    $aI_2+bN$ avec $N=\begin{pmatrix}0&1\\0&0\end{pmatrix}$.
  \item Résoudre alors l'équation, et conclure :
    $X=\pm\begin{pmatrix}1&1/2\\0&1\end{pmatrix}$.
\end{questions}
\end{exercice}

\begin{exercice}[3]{Fibonacci, vu par les matrices}
Soit $A=\begin{pmatrix}1&1\\1&0\end{pmatrix}$ et $(F_n)$ la suite de
Fibonacci ($F_0=0$, $F_1=1$).
\begin{questions}
  \item Démontrer par récurrence que, pour tout $n\geq1$ :
    \[
      A^{n}=\begin{pmatrix}F_{n+1}&F_n\\F_n&F_{n-1}\end{pmatrix}.
    \]
  \item Vérifier que l'application
    $\begin{pmatrix}a&b\\c&d\end{pmatrix}\mapsto ad-bc$ est
    multiplicative sur $\mathcal{M}_2(\K)$ :
    le « $ad-bc$ » d'un produit est le produit des « $ad-bc$ ».
  \item En déduire, en une ligne, l'identité de Cassini :
    $F_{n+1}F_{n-1}-F_n^{2}=(-1)^{n}$ — déjà rencontrée au chapitre 5.
\end{questions}
\end{exercice}

\begin{exercice}[3]{Matrices de permutation}
Pour $\sigma\in S_3$, on note $P_\sigma\in\mathcal{M}_3(\K)$ la matrice
dont la colonne $j$ est la colonne $\sigma(j)$ de\dots{} plus précisément :
$(P_\sigma)_{i,j}=\delta_{i,\sigma(j)}$.
\begin{questions}
  \item Écrire $P_\sigma$ pour la transposition $\tau=(1\,2)$ et pour le
    cycle $c=(1\,2\,3)$.
  \item Montrer que $P_\sigma P_\tau=P_{\sigma\circ\tau}$ : l'application
    $\sigma\mapsto P_\sigma$ est un morphisme de groupes de $S_3$ dans
    $\mathrm{GL}_3(\K)$. Est-il injectif ?
  \item Montrer que $P_\sigma^{-1}=P_\sigma^{\top}$.
\end{questions}
\end{exercice}

\begin{exercice}[3]{Un modèle d'évolution : ville et campagne}
Chaque année, $10\,\%$ des habitants de la ville partent à la campagne, et
$20\,\%$ des habitants de la campagne s'installent en ville. On note $u_n$
et $v_n$ les populations (en milliers) la $n$-ième année, avec
$u_0=v_0=30$.
\begin{questions}
  \item Justifier que $X_{n+1}=AX_n$ avec
    $A=\begin{pmatrix}0{,}9&0{,}2\\0{,}1&0{,}8\end{pmatrix}$, et vérifier
    que la population totale $s_n=u_n+v_n$ est constante.
  \item On pose $d_n=u_n-2v_n$. Montrer que $(d_n)$ est géométrique de
    raison $0{,}7$.
  \item En déduire $u_n$ et $v_n$ en fonction de $n$, puis leurs limites.
  \item Interpréter la limite : vérifier que le vecteur limite est un
    \emph{état d'équilibre}, c'est-à-dire une solution de $AX=X$.
\end{questions}
\end{exercice}

\begin{exercice}[3]{Les complexes, cachés dans $\mathcal{M}_2(\R)$}
Pour $(a,b)\in\R^{2}$, on pose
$M(a,b)=\begin{pmatrix}a&-b\\b&a\end{pmatrix}$, et
$\mathcal{C}=\ens{M(a,b) \mid (a,b)\in\R^{2}}$.
\begin{questions}
  \item Calculer $M(a,b)+M(c,d)$ et $M(a,b)\,M(c,d)$. Que reconnaît-on ?
  \item Montrer que $\mathcal{C}$ est un sous-anneau \emph{commutatif} de
    $\mathcal{M}_2(\R)$, et que l'application
    $a+\mathrm{i}b\mapsto M(a,b)$ est un isomorphisme d'anneaux de $\Cx$
    sur $\mathcal{C}$.
  \item Vérifier que $M(0,1)^{2}=-I_2$ : quelle est l'image de
    $\mathrm{i}$ ?
  \item Montrer que $M(a,b)$ est inversible si et seulement si
    $(a,b)\neq(0,0)$, et faire le lien avec les matrices de rotation :
    $M(\cos\theta,\sin\theta)=R(\theta)$.
\end{questions}
\end{exercice}

% !TEX root = ../../main.tex
%% Défi du chapitre 9 — Carrés magiques et carré de Lo Shu
%% Résultats vérifiés par verifications/defis_06_10.py.

\rubrique{Défi}

\begin{defi}{Carrés magiques et carré de Lo Shu}{1 h 30}
Une matrice $M=(m_{ij})\in\mathcal{M}_3(\R)$ est \emph{magique de somme $s$}
si ses trois lignes, ses trois colonnes et ses deux diagonales ont toutes
pour somme $s$ ; elle est \emph{semi-magique} si l'on ne demande la
condition qu'aux lignes et aux colonnes. On note $J$ la matrice dont tous
les coefficients valent $1$.

\partieExo{A}{Résoudre le système des carrés magiques}
\begin{questions}
  \item Montrer qu'une combinaison linéaire de matrices magiques est
    magique, et que $J$ est magique de somme $3$.
  \item Soit $M$ magique de somme $s$. En additionnant la ligne et la
    colonne du milieu et les deux diagonales, montrer que $m_{22}=s/3$.
  \item Pour $(c,a,b)\in\R^{3}$, on pose
    \[
    L(c,a,b)=\begin{pmatrix}c-a&c+a+b&c-b\\ c+a-b&c&c-a+b\\ c+b&c-a-b&c+a\end{pmatrix}.
    \]
    Vérifier que $L(c,a,b)$ est magique de somme $3c$.
  \item Réciproquement, soit $M$ magique de somme $s$ ; on pose
    $c=m_{22}$, $a=c-m_{11}$ et $b=c-m_{13}$. En utilisant successivement
    les équations du système (diagonales, première ligne, première
    colonne\dots), montrer que $M=L(c,a,b)$. Combien de paramètres
    décrivent les matrices magiques de $\mathcal{M}_3(\R)$ ?
\end{questions}

\partieExo{B}{Le carré de Lo Shu}
\begin{questions}[resume]
  \item On cherche les carrés magiques dont les coefficients sont
    exactement les entiers $1,2,\dots,9$. Montrer que $s=15$ et $c=5$, puis
    que $\ens{|a|,|b|}=\ens{1,3}$. En déduire qu'il y a exactement $8$ tels
    carrés, et que $L(5,1,3)$ est le carré de Lo Shu de la figure.
  \item Vérifier que $L(c,a,b)^{\mathsf T}=L(c,a,-b)$. Expliquer pourquoi les
    huit carrés trouvés sont les images du Lo Shu par les huit symétries du
    carré (défi du chapitre 8).
\end{questions}

\medskip
\begin{center}
\begin{tikzpicture}[line cap=round]
  \foreach \v [count=\k from 0] in {4,9,2,3,5,7,8,1,6}{%
    \pgfmathtruncatemacro{\ii}{div(\k,3)}\pgfmathtruncatemacro{\jj}{mod(\k,3)}
    \ifnum\v=5\fill[otGreen!20] (\jj*1.1,-\ii*1.1) rectangle ++(1.1,-1.1);\fi
    \node[font=\large,text=otInk] at (\jj*1.1+0.55,-\ii*1.1-0.55) {$\v$};}
  \draw[otNavy,line width=0.9pt] (0,0) rectangle (3.3,-3.3);
  \draw[otNavy,line width=0.5pt] (1.1,0) -- (1.1,-3.3) (2.2,0) -- (2.2,-3.3)
    (0,-1.1) -- (3.3,-1.1) (0,-2.2) -- (3.3,-2.2);
  \begin{scope}[xshift=5.2cm]
    \foreach \v [count=\k from 0] in {4,9,2,3,5,7,8,1,6}{%
      \pgfmathtruncatemacro{\ii}{div(\k,3)}\pgfmathtruncatemacro{\jj}{mod(\k,3)}
      \pgfmathsetmacro{\cx}{\jj*1.1+0.55}\pgfmathsetmacro{\cy}{-\ii*1.1-0.55}
      \pgfmathtruncatemacro{\impair}{mod(\v,2)}
      \foreach \d in {1,...,\v}{%
        \pgfmathsetmacro{\rr}{\v>1 ? 0.34 : 0}
        \pgfmathsetmacro{\aa}{90+360*(\d-1)/\v}
        \ifnum\impair=1
          \filldraw[fill=white,draw=otInk,line width=0.5pt] ($(\cx,\cy)+(\aa:\rr)$) circle (2.5pt);
        \else
          \fill[otInk] ($(\cx,\cy)+(\aa:\rr)$) circle (2.5pt);
        \fi}}
    \draw[otNavy!40,line width=0.4pt] (0,0) rectangle (3.3,-3.3);
  \end{scope}
\end{tikzpicture}\\[0.5mm]
\legende{Le carré de Lo Shu, connu en Chine depuis l'Antiquité : chaque
ligne, colonne et diagonale a pour somme $15$. À droite, sa représentation
traditionnelle (impairs en blanc, pairs en noir).}
\end{center}

\partieExo{C}{Calcul matriciel}
\begin{questions}[resume]
  \item Montrer que $M$ est semi-magique de somme $s$ si et seulement si
    $MJ=JM=sJ$. En déduire qu'un produit de matrices semi-magiques de
    sommes $s$ et $t$ est semi-magique de somme $st$, et que l'inverse
    d'une matrice semi-magique inversible de somme $s$ est semi-magique de
    somme $1/s$ (on justifiera que $s\neq0$).
  \item Soit $L=L(5,1,3)$ le Lo Shu. Calculer $L^{2}$ : est-elle
    semi-magique ? magique ?
  \item Vérifier que $L$ est inversible d'inverse
    \[
    L^{-1}=\frac{1}{360}\begin{pmatrix}23&-52&53\\ 38&8&-22\\ -37&68&-7\end{pmatrix},
    \]
    et que $L^{-1}$ est magique : de quelle somme ?
\end{questions}
\end{defi}

